\documentclass[letterpaper,11pt]{article}
\usepackage{resume_style}
\begin{document}

%---------- HEADER ----------
\begin{center}
  {\Huge \scshape Vedaang Chopra} \\ \vspace{2pt}
  \small
  \href{mailto:vedaangchopra@gatech.edu}{\underline{vedaangchopra@gatech.edu}} \;|\;
  \underline{+1 (404) 740-9905} \;|\;
  Atlanta, GA \;|\;
  \href{https://linkedin.com/in/vedaang-chopra}{\underline{LinkedIn}} \;|\;
  \href{https://github.com/Vedaang-Chopra}{\underline{GitHub}} \;|\;
  \href{https://vedaangchopra.live/}{\underline{Website}}
  \\ \vspace{3pt}
  \small
  MS CS (ML) at Georgia Tech \;|\; Advised by \textbf{Dr.\ Matthew Gombolay}, CORE Robotics Lab \\
  \textit{Research:} Agentic computer-use with execution verification -- Multimodal model routing -- Production agentic RAG
\end{center}
\vspace{-6pt}

%---------- EDUCATION ----------
\section{Education}
\resumeSubHeadingListStart
\resumeSubheading
{Georgia Institute of Technology}{Atlanta, GA}
{\textbf{M.S.\ in Computer Science} (Machine Learning), \textbf{GPA: 4.0/4.0}}{Aug 2024 -- Dec 2026}
\vspace{-3pt}
\begin{itemize}[leftmargin=0.20in, topsep=0pt, itemsep=0pt]
  \item\footnotesize{\textbf{Coursework:} Deep Learning, Deep RL, ML Security, Large \& Vision Language Models, Agentic AI, Systems for AI}
\end{itemize}
\vspace{-6pt}
\resumeEducationCompact
{Maharaja Surajmal Institute of Technology}{New Delhi, India}
{B.Tech.\ in Information Technology, CGPA: 8.8/10.0}{Aug 2016 -- Aug 2020}
\resumeSubHeadingListEnd

%---------- RESEARCH EXPERIENCE ----------
\section{Research Experience}
\resumeSubHeadingListStart
\resumeSubheading
{Georgia Tech -- CORE Robotics Lab @ Siemens}{Atlanta, GA}
{Graduate Research Assistant \;|\; Advisor: \textbf{Dr.\ Matthew Gombolay}}{Jan 2026 -- Present}
\resumeItemListStart
  \resumeItem{Building a closed-loop agentic pipeline for parametric CAD code generation: LLM generates CadQuery, execution feedback drives LangGraph repair loop with geometric verification (Chamfer/Hausdorff on STL/B-Rep) -- inference-time search with verifiable reward, 16 modules, 202 tests, VLM visual judgment.}
\resumeItemListEnd

\resumeSubheading
{Georgia Tech -- CS 8903 Special Problems}{Atlanta, GA}
{Graduate Researcher}{Jan 2025 -- Present}
\resumeItemListStart
  \resumeItem{\textbf{ARTEMIS}: Cost-aware VLM routing with trained neural router, SLA-aware load balancing, 6+ backends (Gemma-3, Qwen-VL, DeepSeek-OCR, Llama-4) via vLLM; dynamic accuracy/cost/latency modes.}
  \resumeItem{\textbf{SHASTRA}: Trace-to-graph pipeline converting GAIA agent traces (94 sessions, 4K events) into reusable task-composition graphs; LangGraph planner with Pydantic registry for workflow reuse under cost/latency constraints.}
  \resumeItem{\textbf{ATHENA}: 5-agent supervisor-worker orchestration via LangGraph Command; reflection-based critique, dynamic plan modification, Pydantic structured outputs; cross-modal eval (BLEU-4 +0.18, CLIPScore) on 100 videos.}
\resumeItemListEnd
\resumeSubHeadingListEnd

%---------- INDUSTRY EXPERIENCE ----------
\section{Industry Experience}
\resumeSubHeadingListStart
\resumeSubheading
{Fortinet Technologies Inc.\ (AIOps R\&D)}{Bengaluru, India}
{Software Development Engineer I \& II (ML/AI)}{Feb 2021 -- Jul 2025}
\resumeItemListStart
  \resumeItem{Led agentic RAG diagnostics: LLM plans multi-step tool calls over telemetry, verifies hypotheses via structured function calling, produces root-cause analysis -- reduced resolution time ~70\% (hackathon to production).}
  \resumeItem{Scaled OpenSearch ingestion 50$\to$2,000 events/sec (40x) with Golang; deployed edge ML via ONNX Runtime (~40\% latency reduction); built SLA forecasting for 60+ classifiers.}
  \resumeItem{Patent PCT/IN2022/058026: Unsupervised distributional thresholding for wireless anomaly detection; cut manual troubleshooting ~75\%.}
\resumeItemListEnd
\resumeSubHeadingListEnd

%---------- PUBLICATIONS & PATENTS ----------
\section{Publications \& Patents}
\resumeSubHeadingListStart
\resumeProjectHeading
{\textbf{Patent (Filed):} \href{https://patents.justia.com/inventor/vedaang-chopra}{\textbf{PCT/IN2022/058026}} -- Unsupervised distributional thresholding for wireless anomaly detection.}{}
\resumeSubHeadingListEnd

%---------- TECHNICAL SKILLS ----------
\section{Technical Skills}
\footnotesize
\textbf{Agentic AI:} LangChain, LangGraph, Agentic RAG, Tool/Function Calling, Multi-Agent Orchestration, Supervisor-Worker, Planning \& Reasoning, LLM-as-a-Judge, Structured Generation (Pydantic) \\[1pt]
\textbf{LLMs / VLMs:} vLLM, VLM Routing \& Evaluation, CLIP, SBERT, FAISS, Inference Optimization, Model Evaluation, Cross-Modal Retrieval, Multimodal (video, audio, image, text) \\[1pt]
\textbf{AI / ML:} PyTorch, Hugging Face Transformers, Deep Learning, CNNs, Vision Transformers, NLP, Computer Vision, Scikit-Learn, NumPy, Pandas \\[1pt]
\textbf{Reinforcement Learning:} PPO, DQN, Reward Shaping, Curriculum Learning, Self-Play (academic projects) \\[1pt]
\textbf{Systems:} Python, Go, C/C++, SQL/PostgreSQL, Redis, OpenSearch/Elasticsearch, Docker, Linux, Git, ONNX Runtime, FastAPI, Azure, W\&B, TensorBoard

\end{document}